\documentclass[12pt]{article}
\usepackage[utf8]{inputenc}
\usepackage[portuguese]{babel}
\usepackage{amsmath, amssymb}
\usepackage{graphicx}
\usepackage{float}                  % permite o especificador [H] (posição exata)
\usepackage[margin=2.5cm]{geometry}

\title{1ª Avaliação - Processamento Digital de Sinais}
\author{Ana Iara Loayza Costa \and José Victor Brito Costa \and Milena Freire Britto Neves}
\date{28/07/2026}

\begin{document}
\maketitle

\section*{Questão 1}

\subsection*{a) }

Um sinal discreto é periódico se existir um inteiro $N > 0$ tal que $x[n] = x[n+N]$
para todo $n$. Impondo essa condição a $x[n] = A\cos(\omega_0 n + \theta)$:
\[
A\cos(\omega_0 n + \theta) = A\cos\big(\omega_0(n+N) + \theta\big)
 = A\cos(\omega_0 n + \omega_0 N + \theta),
\]
a igualdade só vale quando $\omega_0 N$ é múltiplo inteiro de $2\pi$. Logo, o sinal é
periódico se, e somente se, existirem $N \in \mathbb{Z}^+$ e $r \in \mathbb{Z}$ tais que:
\[
\omega_0 N = 2\pi r
\]

Para o sinal dado, $\omega_0 = \dfrac{12\pi}{5}$ e $\theta = 0$. Substituindo:
\[
\frac{12\pi}{5}N = 2\pi r \;\Rightarrow\; N = \frac{2\pi r \cdot 5}{12\pi} = \frac{5r}{6}
\]

Como $N$ precisa ser inteiro e $5$ e $6$ são primos entre si, $r$ deve ser múltiplo de
$6$. O menor valor inteiro positivo é $r = 6$, resultando em:
\[
N = \frac{5 \cdot 6}{6} = 5
\]

Portanto, $x[n]$ é periódico, com período fundamental $N = 5$.

\medskip
\noindent\textbf{Verificação direta:}
\[
x[n+5] = \cos\!\left(\frac{12\pi}{5}(n+5)\right)
       = \cos\!\left(\frac{12\pi}{5}n + 12\pi\right)
       = \cos\!\left(\frac{12\pi}{5}n\right) = x[n],
\]
pois $12\pi = 6 \cdot 2\pi$ é múltiplo inteiro de $2\pi$.

\medskip
\noindent\textbf{Frequência equivalente.} Note que $\omega_0 = \frac{12\pi}{5} = 2{,}4\pi$
está fora do intervalo fundamental $-\pi < \omega \le \pi$. Como senoides discretas cujas
frequências diferem de $2\pi$ são indistinguíveis,
\[
\omega_0 = 2{,}4\pi \equiv 2{,}4\pi - 2\pi = 0{,}4\pi,
\]
isto é, $x[n] = \cos(0{,}4\pi n)$, o que leva ao mesmo período:
$0{,}4\pi N = 2\pi r \Rightarrow N = 5r \Rightarrow N = 5$. Essa observação é a chave
para explicar o item (b).

\bigskip
\noindent\textit{Gráfico comprovando a periodicidade:}

\begin{figure}[H]
    \centering
    \includegraphics[width=0.7\textwidth]{images/grafico_1a.png}
\end{figure}

\subsection*{b) }

\textbf{Multiplicando por 2:} $\omega_0' = 2 \cdot \dfrac{12\pi}{5} = \dfrac{24\pi}{5}$
\[
\frac{24\pi}{5}N = 2\pi r \;\Rightarrow\; N = \frac{5r}{12}
\]
$N$ inteiro exige $r$ múltiplo de $12$; o menor é $r = 12 \Rightarrow N = 5$. O período
\textbf{não muda}.

\textit{Por quê?} Usando a frequência equivalente:
$\omega_0' = 4{,}8\pi \equiv 4{,}8\pi - 4\pi = 0{,}8\pi$, ou seja,
$x_2[n] = \cos(0{,}8\pi n)$. Dobrar $\omega_0$ levou a frequência de $0{,}4\pi$ para
$0{,}8\pi$ — ainda dentro de $(0,\pi)$ — e o período continua $5$, pois
$0{,}8\pi N = 2\pi r \Rightarrow N = 2{,}5r \Rightarrow r = 2,\ N = 5$.

\medskip
\textbf{Multiplicando por 5:} $\omega_0'' = 5 \cdot \dfrac{12\pi}{5} = 12\pi$
\[
12\pi N = 2\pi r \;\Rightarrow\; N = \frac{r}{6}
\]
Menor $r$ inteiro válido: $r = 6 \Rightarrow N = 1$.

\textit{Por quê?} $\omega_0'' = 12\pi \equiv 0$: a frequência equivalente é
\textbf{zero}. Como $12\pi n$ é sempre múltiplo inteiro de $2\pi$, temos
$\cos(12\pi n) = 1$ para todo $n$, ou seja, o sinal se torna constante
(período trivial $N=1$).

\medskip
\noindent\textbf{Conclusão.} Nos sinais contínuos, $T = 2\pi/\omega_0$: aumentar a
frequência sempre diminui o período, na proporção inversa. Nos sinais discretos essa
relação não vale — dobrar a frequência manteve $N = 5$ e quintuplicá-la reduziu o sinal
a uma constante. A razão é que apenas a frequência \textbf{módulo $2\pi$} importa, e o
período precisa ser um número \textbf{inteiro} de amostras.

\begin{figure}[H]
    \centering
    \includegraphics[width=0.7\textwidth]{images/grafico_1b_x2.png}
\end{figure}

\begin{figure}[H]
    \centering
    \includegraphics[width=0.7\textwidth]{images/grafico_1b_x5.png}
\end{figure}

\clearpage
\section*{Questão 2}

Foi implementada, em Python, a função \texttt{Deslocamento\_Tempo(x, n, n0)}, que
recebe o sinal $x$, seu vetor de índices $n$ e o escalar de deslocamento $n_0$, e
retorna o vetor $y$ tal que
\[
y[n] = x[n - n_0].
\]

Para cada índice $n_i$ do vetor $n$, a função busca o valor de $x$ correspondente
ao índice $n_i - n_0$ dentro do domínio original de $x$; caso esse índice não
exista nesse domínio, o valor retornado é 0. Dessa forma, $y$ fica definido sobre
o \textbf{mesmo} vetor de índices $n$ recebido como parâmetro, permitindo plotar
$x$ e $y$ diretamente no mesmo eixo, como pedido no enunciado. A implementação é
válida para qualquer vetor $x$, vetor de índices $n$ (não necessariamente
começando em 0 ou contíguo) e escalar $n_0$ (positivo, negativo ou nulo).

\begin{verbatim}
def Deslocamento_Tempo(x, n, n0):
    x = np.asarray(x)
    n = np.asarray(n)
    if x.size != n.size:
        raise ValueError("x e n devem ter o mesmo tamanho.")

    y = np.zeros(n.size, dtype=float)
    for i, ni in enumerate(n):
        alvo = ni - n0
        pos = np.where(n == alvo)[0]
        if pos.size > 0:
            y[i] = x[pos[0]]
    return y
\end{verbatim}

\subsection*{Exemplo}

Como exemplo, considerou-se $x[n]$ um pulso triangular definido para $n \in [-3, 5]$,
com deslocamento $n_0 = 3$ (atraso de 3 amostras). O gráfico abaixo mostra $x[n]$ e
$y[n]$ no mesmo eixo $n$, com legenda identificando cada sinal — note os zeros nas
bordas onde o sinal deslocado ainda não "entrou" no domínio observado.

\begin{figure}[H]
    \centering
    \includegraphics[width=0.7\textwidth]{images/grafico_2.png}
    \caption{Deslocamento no tempo: $x[n]$ e $y[n] = x[n-3]$}
\end{figure}

\clearpage
\section*{Questão 3}

Para um sistema linear e invariante no tempo, a saída é dada pela convolução
entre o sinal de entrada $x[n]$ e a resposta ao impulso $h[n]$:
\[
y[n]=(x*h)[n]=\sum_{k=-\infty}^{\infty}x[k]h[n-k].
\]

Como o enunciado estabelece que ambos os sinais são causais,
$x[n]=h[n]=0$ para $n<0$. Assim, o somatório pode ser limitado a
\[
y[n]=\sum_{k=0}^{n}x[k]h[n-k].
\]

Se $x$ possui $N_x$ amostras e $h$ possui $N_h$ amostras, a convolução completa
possui
\[
N_y=N_x+N_h-1
\]
amostras. Na implementação, cada produto $x[k]h[m]$ contribui para a amostra
de índice $n=k+m$:
\[
y[k+m]\leftarrow y[k+m]+x[k]h[m].
\]
Esse procedimento implementa diretamente a definição, sem usar funções prontas
de convolução.

\begin{verbatim}
def Convolucao(x, h):
    x = np.asarray(x)
    h = np.asarray(h)

    if x.ndim != 1 or h.ndim != 1:
        raise ValueError("x e h devem ser vetores unidimensionais.")
    if x.size == 0 or h.size == 0:
        raise ValueError("x e h não podem ser vetores vazios.")

    tipo = np.result_type(x.dtype, h.dtype)
    y = np.zeros(x.size + h.size - 1, dtype=tipo)

    for k in range(x.size):
        for m in range(h.size):
            y[k + m] += x[k] * h[m]

    return y
\end{verbatim}

\subsection*{Exemplo}

Sejam
\[
x[n]=\{1,2,3\}
\qquad\text{e}\qquad
h[n]=\left\{1,\frac{1}{2}\right\},
\]
com a primeira amostra de cada vetor correspondendo a $n=0$. Então:
\[
\begin{aligned}
y[0]&=x[0]h[0]=1,\\
y[1]&=x[0]h[1]+x[1]h[0]=\frac{1}{2}+2=\frac{5}{2},\\
y[2]&=x[1]h[1]+x[2]h[0]=1+3=4,\\
y[3]&=x[2]h[1]=\frac{3}{2}.
\end{aligned}
\]
Portanto,
\[
\boxed{y[n]=\left\{1,\frac{5}{2},4,\frac{3}{2}\right\}}.
\]

\clearpage
\section*{Questão 4}

A equação fornecida pode ser reescrita como
\[
y[n]=(1+r)y[n-1]+x[n].
\]
Definindo $a=1+r$, obtém-se o sistema de primeira ordem
\[
y[n]-ay[n-1]=x[n].
\]

Como o aporte mensal é constante e o primeiro aporte ocorre ao final do primeiro
mês, sua sequência é $x[n]=x\,u[n-1]$. A resposta ao impulso do sistema é
\[
H(z)=\frac{1}{1-az^{-1}}
\qquad\Longleftrightarrow\qquad
h[n]=a^n u[n].
\]

A resposta total é a soma da resposta à condição inicial com a resposta de estado
nulo. Para um saldo inicial $y[0]=y_0$:
\[
y_{\mathrm{inicial}}[n]=a^n y_0.
\]
A parcela causada pelos aportes é calculada por convolução:
\[
\begin{aligned}
y_{\mathrm{aportes}}[n]
  &=\sum_{k=1}^{n}x\,h[n-k] \\
  &=x\sum_{m=0}^{n-1}a^m.
\end{aligned}
\]
Para $r\neq 0$, a soma geométrica resulta em
\[
y_{\mathrm{aportes}}[n]
=x\frac{a^n-1}{a-1}
=x\frac{(1+r)^n-1}{r}.
\]
Portanto, a solução fechada, sem iteração recursiva, é
\[
\boxed{
y[n]=y_0(1+r)^n+x\frac{(1+r)^n-1}{r}
},\qquad r\neq 0.
\]
Quando $r=0$:
\[
\boxed{y[n]=y_0+nx}.
\]

\begin{verbatim}
def Investimento(y0, r, x, n):
    if n < 0:
        raise ValueError("n nao pode ser negativo.")
    if r <= -1:
        raise ValueError("r deve ser maior que -1.")

    if r == 0:
        return y0 + n*x

    fator = (1 + r)**n                  # (1+r)^n
    return y0*fator + x*(fator - 1)/r   # y0(1+r)^n + x[(1+r)^n - 1]/r
\end{verbatim}

\subsection*{Exemplo}

Para $y_0=1000$, rentabilidade mensal de $1\%$ ($r=0{,}01$), aporte mensal de
$200$ e prazo de $12$ meses:
\[
\begin{aligned}
\mathrm{Investimento}(1000,0{,}01,200,12)
&=1000(1{,}01)^{12}
  +200\frac{(1{,}01)^{12}-1}{0{,}01}\\
&\approx 3663{,}33.
\end{aligned}
\]

\clearpage
\section*{Questão 5}

Dada a equação de diferenças
\[
y[n] = 0{,}038\,x[n-1] + 0{,}034\,x[n-2] + 1{,}63\,y[n-1] - 0{,}70\,y[n-2],
\]
determinam-se a seguir a função de transferência, a resposta em frequência e a
saída para uma entrada senoidal.

\subsection*{a) Função de transferência $H(z)$}

\noindent\textbf{Usando a forma padrão.} Passando os termos de saída para o lado
esquerdo:
\[
y[n] - 1{,}63\,y[n-1] + 0{,}70\,y[n-2] = 0{,}038\,x[n-1] + 0{,}034\,x[n-2].
\]

\noindent\textbf{Agora temos a transformada $z$.} Pela propriedade do deslocamento,
$x[n-k] \leftrightarrow z^{-k}X(z)$ (válida com condições iniciais nulas, hipótese
usada para definir a função de transferência):
\[
Y(z) - 1{,}63\,z^{-1}Y(z) + 0{,}70\,z^{-2}Y(z)
 = 0{,}038\,z^{-1}X(z) + 0{,}034\,z^{-2}X(z).
\]

\noindent\textbf{Evidenciando $Y(z)$ e $X(z)$:}
\[
Y(z)\left(1 - 1{,}63\,z^{-1} + 0{,}70\,z^{-2}\right)
= X(z)\left(0{,}038\,z^{-1} + 0{,}034\,z^{-2}\right).
\]

\noindent\textbf{Isolando $Y(z)/X(z)$:}
\[
\boxed{
H(z) = \frac{Y(z)}{X(z)}
= \frac{0{,}038\,z^{-1} + 0{,}034\,z^{-2}}{1 - 1{,}63\,z^{-1} + 0{,}70\,z^{-2}}
}
\]
Multiplicando numerador e denominador por $z^{2}$, obtém-se a forma em potências
positivas:
\[
H(z) = \frac{0{,}038\,z + 0{,}034}{z^{2} - 1{,}63\,z + 0{,}70}.
\]

\noindent\textbf{Polos e zeros.} O zero finito vem de $0{,}038z + 0{,}034 = 0$, isto é,
$z = -0{,}034/0{,}038 = -0{,}895$. Os polos vêm de $z^{2} - 1{,}63z + 0{,}70 = 0$:
\[
z = \frac{1{,}63 \pm \sqrt{1{,}63^{2} - 4(0{,}70)}}{2}
  = \frac{1{,}63 \pm \sqrt{-0{,}1431}}{2}
  = 0{,}815 \pm 0{,}189j.
\]
Como o produto das raízes é o termo independente, $z_1 z_2 = |z|^2 = 0{,}70$, logo
\[
|z| = \sqrt{0{,}70} \approx 0{,}8367 < 1.
\]

\noindent\textbf{Região de convergência e estabilidade.} O sistema é causal (a saída
depende apenas de amostras passadas), então a RDC é $|z| > 0{,}8367$, o exterior do
círculo que passa pelos polos. Como essa RDC \textbf{contém o círculo unitário}
$|z| = 1$, o sistema é BIBO-estável e a resposta em frequência $H(e^{j\omega})$
existe --- condição necessária para os itens seguintes.

\subsection*{b) Representação de $H(z)$ em função da frequência}

A resposta em frequência é obtida avaliando $H(z)$ sobre o círculo unitário, isto é,
fazendo $z = e^{j\omega}$ (passo legítimo porque a RDC contém $|z| = 1$):
\[
H(e^{j\omega}) =
\frac{0{,}038\,e^{-j\omega} + 0{,}034\,e^{-j2\omega}}
     {1 - 1{,}63\,e^{-j\omega} + 0{,}70\,e^{-j2\omega}}
= \frac{N(e^{j\omega})}{D(e^{j\omega})}.
\]

\noindent\textbf{Separando parte real e imaginária.} Aplicando a fórmula de Euler,
$e^{-jk\omega} = \cos(k\omega) - j\,\mathrm{sen}(k\omega)$:
\[
\begin{aligned}
N(e^{j\omega}) &= \underbrace{\big[0{,}038\cos\omega + 0{,}034\cos 2\omega\big]}_{N_R(\omega)}
 + j\underbrace{\big[-0{,}038\,\mathrm{sen}\,\omega - 0{,}034\,\mathrm{sen}\,2\omega\big]}_{N_I(\omega)},\\[2mm]
D(e^{j\omega}) &= \underbrace{\big[1 - 1{,}63\cos\omega + 0{,}70\cos 2\omega\big]}_{D_R(\omega)}
 + j\underbrace{\big[1{,}63\,\mathrm{sen}\,\omega - 0{,}70\,\mathrm{sen}\,2\omega\big]}_{D_I(\omega)}.
\end{aligned}
\]

Como o módulo de um quociente é o quociente dos módulos, e a fase de um quociente é a
diferença das fases:
\[
\boxed{
\left|H(e^{j\omega})\right| =
\frac{\sqrt{N_R^2(\omega) + N_I^2(\omega)}}{\sqrt{D_R^2(\omega) + D_I^2(\omega)}}
}
\qquad
\boxed{
\angle H(e^{j\omega}) =
\operatorname{arctg}\frac{N_I(\omega)}{N_R(\omega)}
- \operatorname{arctg}\frac{D_I(\omega)}{D_R(\omega)}
}
\]
Nas contas numéricas usa-se \texttt{arctan2}, que resolve corretamente o quadrante.

\medskip
\noindent\textbf{Forma compacta da magnitude via $|H|^2 = HH^{*}$.} Colocando
$e^{-j\omega}$ em evidência no numerador,
$N = e^{-j\omega}\big(0{,}038 + 0{,}034\,e^{-j\omega}\big)$, e como
$|e^{-j\omega}| = 1$:
\[
\begin{aligned}
|N|^2 &= \left|0{,}038 + 0{,}034e^{-j\omega}\right|^2
= 0{,}038^2 + 0{,}034^2 + 2(0{,}038)(0{,}034)\cos\omega\\
      &= 0{,}002600 + 0{,}002584\cos\omega.
\end{aligned}
\]
Fazendo $|D|^2 = DD^{*}$ e usando $e^{-jk\omega} + e^{jk\omega} = 2\cos k\omega$:
\[
\begin{aligned}
|D|^2 &= \big(1 + 1{,}63^2 + 0{,}70^2\big)
 - 2\big(1{,}63 + 1{,}63 \cdot 0{,}70\big)\cos\omega + 2(0{,}70)\cos 2\omega\\
      &= 4{,}1469 - 5{,}542\cos\omega + 1{,}40\cos 2\omega.
\end{aligned}
\]
Portanto:
\[
\boxed{
\left|H(e^{j\omega})\right| =
\sqrt{\frac{0{,}002600 + 0{,}002584\cos\omega}
           {4{,}1469 - 5{,}542\cos\omega + 1{,}40\cos 2\omega}}
}
\]
As duas formas são equivalentes --- a verificação numérica feita no notebook acusa
erro máximo da ordem de $10^{-14}$.

\medskip
\noindent\textbf{Intervalo de $\omega$.} Como os coeficientes da equação de diferenças
são reais, vale a simetria hermitiana $H(e^{-j\omega}) = H^{*}(e^{j\omega})$: a
magnitude é \textbf{par} e a fase é \textbf{ímpar}. Além disso, $H(e^{j\omega})$ é
periódica de período $2\pi$. Por isso basta representar $\omega$ no intervalo
$-\pi < \omega \le \pi$.

\subsection*{c) Gráficos da magnitude e da fase}

A Figura~\ref{fig:magfase} mostra a magnitude e a fase de $H(e^{j\omega})$ no
intervalo fundamental $-\pi < \omega \le \pi$. Confirma-se a simetria prevista no item
(b): magnitude par e fase ímpar.

O sistema comporta-se como um \textbf{filtro passa-baixas}: o ganho vale
$|H(e^{j0})| = 1{,}0286$ em frequência nula, apresenta um leve pico de ressonância de
$1{,}0589$ em $\omega \approx 0{,}14$ rad (efeito dos polos complexos, situados no
ângulo $\pm 0{,}228$ rad) e decai a $|H(e^{j\pi})| = 0{,}0012$ na frequência de
Nyquist. A frequência de corte de $-3$\,dB fica em torno de $\omega \approx 0{,}33$
rad ($\approx 0{,}10\pi$), de modo que $\omega = \pi/4$ (linha pontilhada) já está bem
dentro da banda de rejeição.

Na curva de fase são traçadas duas versões: o \textbf{valor principal}, devolvido por
\texttt{arctan2} e sempre restrito a $(-\pi, \pi]$, e a \textbf{fase desenrolada}
(\textit{unwrap}). Os saltos verticais de $2\pi$ da primeira curva são um artefato da
representação --- e não descontinuidades do sistema ---, o que fica evidente ao
compará-la com a segunda, contínua.

\begin{figure}[H]
    \centering
    \includegraphics[width=0.85\textwidth]{images/grafico_5_mag_fase.png}
    \caption{Magnitude e fase de $H(e^{j\omega})$ em $-\pi < \omega \le \pi$.}
    \label{fig:magfase}
\end{figure}

\subsection*{d) Saída para $x[n] = \cos\!\left(\frac{\pi}{4}n\right)$}

\noindent\textbf{Propriedade utilizada.} Para um sistema LIT estável, a exponencial
complexa é autofunção: se a entrada é $e^{j\omega_0 n}$, a saída é
$H(e^{j\omega_0})e^{j\omega_0 n}$. Escrevendo o cosseno como soma de exponenciais e
usando a simetria hermitiana, obtém-se o resultado clássico: para
$x[n] = \cos(\omega_0 n)$,
\[
y[n] = \left|H(e^{j\omega_0})\right|\,
\cos\!\big(\omega_0 n + \angle H(e^{j\omega_0})\big).
\]
Ou seja, o sistema \textbf{não cria novas frequências}: apenas escala a amplitude por
$|H|$ e desloca a fase por $\angle H$.

\medskip
\noindent\textbf{Inicialmente, identifica-se a frequência.} Comparando
$x[n] = \cos(\pi n/4)$ com $\cos(\omega_0 n)$: $\omega_0 = \dfrac{\pi}{4}$.

\noindent\textbf{Avaliando os exponenciais.}
\[
e^{-j\pi/4} = \cos\tfrac{\pi}{4} - j\,\mathrm{sen}\tfrac{\pi}{4}
 = 0{,}70711 - 0{,}70711j,
\qquad
e^{-j\pi/2} = -j.
\]

\noindent\textbf{Numerador e denominador:}
\[
\begin{aligned}
N &= 0{,}038(0{,}70711 - 0{,}70711j) + 0{,}034(-j) = 0{,}026870 - 0{,}060870j,\\
D &= 1 - 1{,}63(0{,}70711 - 0{,}70711j) + 0{,}70(-j)\\
  &= 1 - 1{,}152584 + 1{,}152584j - 0{,}70j = -0{,}152584 + 0{,}452584j.
\end{aligned}
\]

\noindent\textbf{Módulos e fases:}
\[
\begin{aligned}
|N| &= \sqrt{0{,}026870^2 + 0{,}060870^2} = 0{,}066537,\\
\angle N &= \operatorname{arctg2}(-0{,}060870;\ 0{,}026870) = -1{,}15509 \text{ rad},\\[1mm]
|D| &= \sqrt{0{,}152584^2 + 0{,}452584^2} = 0{,}477613,\\
\angle D &= \operatorname{arctg2}(0{,}452584;\ -0{,}152584) = +1{,}89597 \text{ rad}.
\end{aligned}
\]

\noindent\textbf{Combinando:}
\[
\left|H(e^{j\pi/4})\right| = \frac{|N|}{|D|}
 = \frac{0{,}066537}{0{,}477613} = 0{,}13931,
\]
\[
\angle H(e^{j\pi/4}) = \angle N - \angle D
 = -1{,}15509 - 1{,}89597 = -3{,}05106 \text{ rad} \approx -174{,}81^\circ.
\]

O mesmo valor de magnitude resulta da fórmula fechada do item (b):
$|N|^2 = 0{,}002600 + 0{,}002584\cos\frac{\pi}{4} = 0{,}0044272$ e
$|D|^2 = 4{,}1469 - 5{,}542\cos\frac{\pi}{4} + 1{,}40\cos\frac{\pi}{2} = 0{,}2281142$,
de onde $|H| = \sqrt{0{,}0044272/0{,}2281142} = 0{,}13931$.

\medskip
\noindent\textbf{Resultado:}
\[
\boxed{
y[n] = 0{,}13931\,\cos\!\left(\frac{\pi}{4}n - 3{,}05106\right)
}
\]

\noindent\textbf{Interpretação.} Como $\omega_0 = \pi/4$ está na banda de rejeição do
filtro passa-baixas, a senoide é atenuada a cerca de $14\%$ da amplitude original. A
fase de $\approx -3{,}05$ rad (quase $-\pi$) significa que a saída fica praticamente
\textbf{invertida} em relação à entrada.

\medskip
\noindent\textbf{Observação sobre o regime.} A expressão acima é a resposta ao cosseno
\textit{permanente} (aplicado desde $n \to -\infty$), convenção adotada em sala. Se a
senoide fosse ligada apenas em $n = 0$, haveria ainda um transitório devido aos polos,
que decai como $|z|^n = 0{,}8367^n$ e se torna desprezível após algumas dezenas de
amostras.

\subsection*{e) Gráfico de $x[n]$ e $y[n]$}

A Figura~\ref{fig:xy} apresenta a entrada $x[n]$ e a saída $y[n]$ no mesmo eixo, com
legendas discriminando os sinais. Observa-se a forte redução de amplitude
($1 \rightarrow 0{,}139$) e a inversão de fase previstas analiticamente.

Para validar a dedução do item (d), o notebook compara essa saída analítica com a
obtida \textbf{iterando diretamente a equação de diferenças original}. As duas
coincidem com erro máximo da ordem de $10^{-16}$, ou seja, dentro da precisão numérica
da máquina --- o que confirma tanto a $H(z)$ obtida quanto a expressão final de $y[n]$.

\begin{figure}[H]
    \centering
    \includegraphics[width=0.85\textwidth]{images/grafico_5_xy.png}
    \caption{Entrada $x[n]=\cos(\pi n/4)$ e saída $y[n]$ do sistema.}
    \label{fig:xy}
\end{figure}

\end{document}
